\section{Introdução e Repositórios}
O objetivo desta etapa do projeto de PDI consiste em aumentar a variabilidade e a quantidade de imagens da base original em 3 vezes, aplicando transformações de intensidade e filtragem espacial linear. O ambiente de execução e armazenamento do projeto utiliza os seguintes endereços oficiais:
\begin{itemize}
    \item \textbf{Repositório Oficial no GitHub:} \\
    \url{https://github.com/HayJM/Projeto.git}
    \item \textbf{Jupyter Notebook Executável (Google Colab):} \\
    \url{https://colab.research.google.com/drive/1P4VxhGKD7LMfsBAVPdxCGIpQoDJSre4r?usp=sharing}
\end{itemize}

\section{Formalização Matemática das Transformações}

\subsection{Transformação Logarítmica}
A transformação logarítmica é utilizada para a expansão de contraste, mapeando uma gama estreita de níveis de intensidade baixa (tons escuros) em uma gama mais ampla de níveis de saída. A expressão matemática que rege esta transformação é:
\begin{equation}
s = c \cdot \log(1 + r)
\end{equation}
Onde $r$ é a intensidade do píxel original ($r \geq 0$) e $c$ é uma constante de escala calculada dinamicamente para garantir o mapeamento do valor máximo de saída no limite de 8 bits (255):
\begin{equation}
c = \frac{255}{\log(1 + \max(f))}
\end{equation}

\subsection{Transformação Exponencial (Lei da Potência / Correção Gamma)}
A transformação exponencial permite expandir ou comprimir o contraste de forma não linear através de uma curva controlada pelo parâmetro $\gamma$ (gamma):
\begin{equation}
s = c \cdot r^{\gamma}
\end{equation}
Onde $r$ representa a intensidade do píxel original previamente normalizada para o intervalo $[0, 1]$. Para o algoritmo implementado, adotou-se o expoente de controle $\gamma = 0.5$ para expandir os tons escuros.

\subsection{Filtro da Média usando Convolução}
O filtro da média é uma operação de filtragem espacial linear de suavização (passa-baixo), que substitui o valor de cada píxel pela média aritmética das intensidades contidas em uma vizinhança retangular definida por um \textit{kernel} (máscara). A operação de convolução discreta bidimensional para uma vizinhança de tamanho $m \times n$ é dada por:
\begin{equation}
g(x, y) = \sum_{s=-a}^{a} \sum_{t=-b}^{b} w(s, t) \cdot f(x+s, y+t)
\end{equation}

Para o Filtro da Média com Kernel $5 \times 5$, todos os 25 coeficientes possuem pesos idênticos e normalizados, conforme representado matricialmente de forma completa:
\begin{equation}
w = \frac{1}{25} \begin{bmatrix}
1 & 1 & 1 & 1 & 1 \\
1 & 1 & 1 & 1 & 1 \\
1 & 1 & 1 & 1 & 1 \\
1 & 1 & 1 & 1 & 1 \\
1 & 1 & 1 & 1 & 1
\end{bmatrix}
\end{equation}

A equação simplificada aplicada a cada píxel do mapa de intensidade resulta em:
\begin{equation}
g(x, y) = \frac{1}{25} \sum_{i=-2}^{2} \sum_{j=-2}^{2} f(x+i, y+j)
\end{equation}

\section{Estrutura de Execução e Classes}
O bloco de código implementado detecta e processa dinamicamente as 10 classes da base de dados coletada:
\texttt{00\_Cherry\_Tomatoes}, \texttt{01\_Fig}, \texttt{02\_Gala\_Apple}, \texttt{03\_Pear}, \texttt{04\_Persimmon}, \texttt{05\_Pitanga}, \texttt{06\_Pomegranate}, \texttt{07\_Sicilian\_Lemon}, \texttt{08\_Strawberry} e \texttt{09\_Tangerine}.

A rotina realiza o clone automático do repositório, efetua as transformações matemáticas citadas e conserva os resultados gerando a estrutura final nomeada como \texttt{augmented\_dataset}.

